\section{Conclusion}
\label{sec:conclusion}

This paper evaluated whether the elementary mechanisms of a compound motor
fault can be recovered when their exact combination, their compound context,
and the operating condition are varied independently of the training history.
On a public multimodal dataset of an induction motor whose compound faults
span mechanical and electrical modalities, a locked linear constituent probe
over physical features specific to each mechanism showed that strong closed
set discriminability (accuracy 0.913 over the 24 diagnostic states, 95\% CI
$[0.885, 0.938]$) coexists with exact constituent set recovery of only 0.083
to 0.157 on unseen compounds. The two results measure different abilities,
and the gap between them is the subject of this paper.

Four conclusions follow from the controlled analyses. First, at the aggregate
level, the dominant failure mode is incomplete decomposition rather than
complete diagnostic silence. Constituent recall remained between 0.630 and
0.694, but a constant predictor driven by prevalence alone could itself reach
0.667 recall at the cost of 1.667 false additions per run and zero exact
matches, so constituent recall must be read jointly with false additions,
micro averaged $F_1$, and exact set recovery rather than on its own. The constituent resolved controls
further showed that these failures do not share a common origin. Dynamic
eccentricity was almost perfectly learnable as an isolated fault (AUROC
0.9995 and $F_1 = 0.923$) but collapsed to compound recall of 0.017, which is
the clearest evidence that poor compound recovery cannot be explained by weak
detectability in isolation alone. Static eccentricity remained highly
learnable in isolation and rose from 0.577 to 0.866 in $F_1$ under the RBF
SVM, indicating pronounced sensitivity to classifier family, whereas winding
was already less specific in isolation and was rescued by none of the
nonlinear alternatives. Bearing inner, bearing outer, and broken bar evidence
remained strongly discriminative throughout. Constituent level failure in
this dataset therefore exhibits distinct empirical patterns rather than a
uniform cross modal deficit.

Second, the adverse effect of compound context deprivation was larger when
the operating condition was simultaneously unseen in the primary analysis.
Clustering at the composition level supported the interaction for constituent
recall ($+0.051$, 95\% CI $[0.009, 0.093]$) and for recovery of all
constituents ($+0.083$, $[0.009, 0.157]$), while the interaction in micro
averaged $F_1$ remained unresolved ($+0.040$, $[-0.004, 0.086]$). After
matching the recombination training set cardinality to the two sided regime,
all three interactions remained positive on average but none was resolved.
The observed context effect is therefore sensitive to the available training
support and should not be interpreted as an effect of compound context
exposure that is independent of training cardinality.

Third, constituent evidence transfers across compound partners often but
directionally. Detectability within a pair was near perfect throughout, yet
transfer across partners ranged from perfect to a statistically supported
reversal of score orientation (Holm adjusted $p \le 0.014$), and the
separability of distinct recordings of the same fault state requires that
reversal to be read as context dependent score transfer rather than proven
inversion of a physical signature. Fourth, absolute
decomposition quality is a property of the operating point and of the model.
Thresholding weighted toward precision raised exact match recovery to 0.157
to 0.287 at reduced recall, and nonlinear classifier families
reduced false additions from approximately 0.68 to approximately 0.26 per
run, while the recombination advantage over two sided withholding under
unseen conditions remained directionally positive across all four classifier
families.

For practice, these results caution against reading closed set accuracy as
deployment readiness and argue for calibrating constituent detectors with an
explicit cost structure over precision and recall. They are also consistent
with compound context novelty and operating condition novelty being more
demanding in combination than either alone, but because that interaction did
not survive matching of training cardinality, this is offered as a direction
for commissioning practice to test rather than as an established requirement.
Finally, the dataset provides insufficient replicated recordings of the same
state to separate fault, assembly, sensor mounting, thermal, and recording
related effects, so the transfer analyses stop at carefully bounded
attribution. Representation strategies that promote consistency of
constituent scores across partners remain hypotheses rather than established
remedies. The protocol, the prevalence references, and the controls used here
are directly portable to the companion gearbox release and to other
multimodal datasets.
